\chapter*{Introduzione}
\chaptermark{Introduzione}

La determinazione della natura della materia oscura costituisce uno dei problemi aperti più profondi della fisica fondamentale contemporanea. Sebbene molteplici evidenze sperimentali indipendenti su scala astrofisica e cosmologica (dalle anomalie cinematiche nelle curve di rotazione galattiche al comportamento delle galassie in collisione come lente gravitazionale) convergano nell'attestare che oltre l'80\% della materia presente nell'Universo non sia di natura barionica, la sua composizione rimane a oggi sconosciuta.

I candidati teorici più accreditati sono le WIMP (\textit{Weakly Interacting Massive Particles}), particelle massive neutre che nell'Universo primordiale si creavano e annichilivano continuamente fino al momento del freeze-out, in cui il progressivo raffreddamento ha arrestato prima i processi di creazione e poi quelli di annichilazione. Per riprodurre l'abbondanza cosmologica osservata oggi, la sezione d'urto di annichilazione richiesta deve essere tipica dell'interazione nucleare debole: questa notevole coincidenza teorica (nota come WIMP miracle) rende tale candidato particolarmente solido e motiva la progettazione di esperimenti dedicati all'osservazione diretta di questa interazione. Negli ultimi decenni, imponenti esperimenti di rivelazione diretta basati su bersagli a gas nobili liquefatti hanno esplorato estesamente la finestra standard delle WIMP con masse dell'ordine dei $\text{GeV} - \text{TeV}$, escludendo porzioni considerevoli dello spazio dei parametri. Oggi, la comunità scientifica rivolge un crescente interesse verso candidati di materia oscura leggeri, con masse sub-$\text{GeV}$. In questa finestra cinematica, gli urti elastici contro i nuclei bersaglio rilasciano energie di rinculo estremamente esigue, inferiori al keV, risultando inaccessibili ai rivelatori tradizionali ad alta soglia e richiedendo tecnologie di sensoristica innovative ad altissima sensibilità.

In tale scenario si colloca l'esperimento CRESST (\textit{Cryogenic Rare Event Search with Superconducting Thermometers}), operativo presso i Laboratori Nazionali del Gran Sasso (LNGS). CRESST impiega cristalli scintillanti di ($\text{CaWO}_4$) raffreddati a circa $10\,\text{mK}$ mediante criostati a diluizione $^3\text{He}/^4\text{He}$. La lettura combinata del calore (tramite sensori a transizione di fase, TES, e amplificatori SQUID) e della contemporanea luce di scintillazione consente sia una straordinaria discriminazione del fondo radioattivo sia il raggiungimento di soglie energetiche dell'ordine di poche decine di elettronvolt, posizionando l'esperimento all'avanguardia mondiale nell'indagine della materia oscura leggera.

Il raggiungimento di soglie così basse e la corretta interpretazione di eventuali segnali fisici pongono due sfide: il trattamento del rumore sperimentale e la determinazione accurata della scala energetica del calorimetro. Il presente lavoro di tesi affronta sistematicamente entrambe queste problematiche attraverso l'analisi dei dati acquisiti dal rivelatore \texttt{TUM93A}.

L'elaborato è organizzato in cinque capitoli:

\begin{itemize}
    \item Il \textbf{Capitolo 1} delinea il quadro teorico e fenomenologico della materia oscura, ripercorrendo le principali evidenze sperimentali, le classi di candidati particellari (WIMP e assioni) e lo stato dell'arte delle tecniche di rivelazione diretta, indiretta.
    
    \item Il \textbf{Capitolo 2} presenta l'apparato sperimentale di CRESST, approfondendo il principio della lettura a doppio canale (calore e luce) e la discriminazione del fondo tramite \textit{Light Yield}. Vengono inoltre descritti i sensori superconduttori TES accoppiati a SQUID, l'articolata schermatura passiva e attiva contro i muoni cosmici e il funzionamento del criostato a diluizione.

    \item Il \textbf{Capitolo 3} è dedicato alla caratterizzazione del rumore sperimentale e allo sviluppo dell'algoritmo di filtraggio digitale. Viene formulato l'Optimum Filter (OF) nel dominio delle frequenze a partire dalla stima della densità spettrale di potenza del rumore (\textit{Noise Power Spectrum}) e dai template per impulsi di particella e impulsi termici artificiali (\textit{testpulses}). L'applicazione del filtro permette di quantificare l'efficienza di rivelazione mediante simulazioni e di definire una soglia operativa di trigger al livello dell'1\% di eventi spuri, corrispondente a una soglia energetica equivalente di circa $37\,\text{eV}$.

    \item Il \textbf{Capitolo 4} affronta la calibrazione energetica del calorimetro attraverso l'analisi congiunta di una sorgente interna a bassa energia di $^{55}\text{Fe}$ e di una sorgente $\gamma$ penetrante ad alta energia di $^{57}\text{Co}$. Viene trattata la risposta non lineare del sensore e la ricostruzione degli impulsi saturi tramite algoritmi a template troncato (\texttt{cait}). La validità dell'estrapolazione della curva di calibrazione verso la regione sub-keV viene confermata dal confronto con le emissioni del $^{55}\text{Fe}$ (entro uno scarto relativo inferiore al $5\%$) e applicata infine alla predizione dell'ampiezza attesa per il rinculo nucleare monoenergetico a $112.4\,\text{eV}$ indotto da cattura di neutroni termici su $^{182}\text{W}$.

    \item Il \textbf{Capitolo 5} raccoglie le conclusioni del lavoro discutendo le prospettive applicative per le future campagne di misura di CRESST-III.
\end{itemize}